\subsectionaddtolist{ شباهت‌های بین دو الگوی \lr{Clean Architecture} و \lr{MVVM} }
با توجه به مباحث مطرح‌شده در متن مقاله و \lr{Case Study}، شباهت‌های بنیادی و ساختاری زیر بین این دو الگو دیده می‌شود:

\begin{itemize}
    \item \textbf{\lr{Separation of Concerns}:} هر دو الگو با هدف تفکیک وظایف لایه‌های مختلف سیستم طراحی شده‌اند. الگوی \lr{MVVM} رابط کاربری را از منطق نمایش جدا می‌کند و \lr{Clean Architecture} کل نرم‌افزار را به لایه‌های مستقل تقسیم می‌نماید.
    \item \textbf{ \lr{Testability}:} در هر دو ساختار، بستری عالی برای \lr{Unit Testing} مهیا شده است. در \lr{MVVM}، لایه \lr{ViewModel} مستقل از \lr{UI} تست می‌شود و در \lr{Clean Architecture} نیز هر لایه به طور جداگانه قابل ارزیابی است.
    \item \textbf{ \lr{Modularity}:} هر دو الگو کدها را به بخش‌های مجزا و مستقل تقسیم می‌کنند که توسعه موازی و مدیریت پروژه را تسهیل می‌کند.
    \item \textbf{پیچیدگی و سربار در پروژه‌های کوچک:} مقاله به صراحت اشاره می‌کند که هر دو الگو برای پروژه‌های ساده و کوچک، به دلیل ایجاد کدهای تکراری زیاد و پیچیدگی اولیه، نوعی زیاده‌روی به شمار می‌روند .
\end{itemize}

% ----------------------------------------------------------------------
% سوال دوم
% ----------------------------------------------------------------------
\subsectionaddtolist{ نحوه ادغام و ترکیب دو الگو توسط نویسنده}
نویسنده مقاله برای بهره‌گیری از نقاط قوت هر دو ساختار، آن‌ها را به شکل یک معماری هیبریدی منسجم ادغام کرده است. نقطه عطف این ترکیب در لایه واسط قرار دارد.
نویسنده \lr{ViewModel} (از الگوی \lr{MVVM}) را دقیقاً هم‌تراز و معادل لایه \lr{Interface Adapter} (از الگوی \lr{Clean Architecture}) قرار داده است با این کار، \lr{ViewModel} وظیفه مدیریت وضعیت نمایش را بر عهده می‌گیرد در حالی که با لایه‌های درونی‌تر برای اجرای منطق تجاری در ارتباط است.

{جریان ارتباطی لایه‌ها در لایه‌بندی ترکیبی:}
\begin{enumerate}
    \item \textbf{\lr{View} (لایه \lr{UI}):} بیرونی‌ترین سطح که با کاربر در تعامل است و از طریق فرآیند \lr{Data Binding} به ویژگی‌های \lr{ViewModel} متصل می‌شود.
    \item \textbf{\lr{ViewModel} (لایه \lr{Presentation}):} نقش واسط را ایفا کرده، ورودی‌های \lr{View} را گرفته و \lr{Use Case} متناسب را فراخوانی می‌کند.
    \item \textbf{\lr{Use Cases} (لایه \lr{Business Logic}):} حاوی منطق‌های خاص هر قابلیت برنامه (مانند افزودن یا حذف یک \lr{Task}) است.
    \item \textbf{\lr{Entities} (لایه \lr{Core Rules}):} شامل اشیاء پایه و قوانین بنیادی کسب‌وکار است که به هیچ فریمورک خارجی وابستگی ندارد.
    \item \textbf{\lr{Data Layer}:} وظیفه تعامل با پایگاه‌داده محلی یا \lr{API}های سرور و مدیریت ماندگاری داده‌ها را بر عهده دارد.
\end{enumerate}

% ----------------------------------------------------------------------
% سوال سوم
% ----------------------------------------------------------------------
\subsectionaddtolist{ پروژه‌های مناسب برای استفاده از این الگوی ترکیبی و دلایل آن}
بر اساس تحلیل مزایا و معایب در متن مقاله، این معماری ترکیبی برای پروژه‌هایی با ویژگی‌های زیر بیشترین بازدهی را دارد:

\begin{itemize}
    \item \textbf{پروژه‌های بزرگ و سازمانی:} به دلیل ساختار کاملاً ماژولار، این الگو اجازه می‌دهد تیم‌های بزرگ بدون تداخل با یکدیگر روی بخش‌های مختلف پروژه کار کنند.
    \item \textbf{پروژه‌های با طول عمر بالا و نیازمند نگهداری:} از آنجا که لایه‌ها به وسیله اینترفیس‌ها از هم ایزوله شده‌اند، تغییر در یک لایه (مثلاً تعویض کتابخانه شبکه یا تغییر دیتابیس) تأثیری روی منطق اصلی برنامه نخواهد گذاشت.
    \item \textbf{پروژه‌های مبتنی بر \lr{TDD}:} برای برنامه‌هایی که پوشش تست بالایی نیاز دارند، تفکیک دقیق لایه‌ها امکان نوشتن تست‌های واحدِ پایدار و مستقل را فراهم می‌سازد.
\end{itemize}

پیاده‌سازی این ساختار نیازمند تعریف اینترفیس‌های متعدد، کلاس‌های مختلف و تنظیمات اولیه پیچیده‌ای است در پروژه‌های کوچک، زمان و هزینه‌ای که صرف نوشتن این کدهای اضافه می‌شود، توجیه اقتصادی و فنی ندارد.